% !TeX root = ../../Book4.tex
% 纲要标题：### 1.5 范畴代数与图表表示
% 纲要细目（写作范围）：
% ---
% #### 1.5.1 复合的线性化与范畴代数
% #### 1.5.2 图表与左模的对应
% #### 1.5.3 群、路径与偏序的例子

\section{范畴代数与图表表示}
\label{b4:sec:0105}

给出两个 \(k\)-模 \(V_0,V_1\) 及一个线性映射 \(u:V_0\to V_1\)，便得到最简单的模图表。令 \(V=V_0\oplus V_1\)，记投影为 \(p_i:V\rightarrow V_i\)，包含为 \(\iota_i:V_i\rightarrow V\)。作用在总模 \(V\) 上的是幂等自同态 \(e_i=\iota_i\circ p_i\)，它取出第 \(i\) 个分量再放回原位；映射 \(u\) 则给出自同态 \(\iota_1\circ u\circ p_0\)，把第 \(0\) 个分量送入第 \(1\) 个分量，并在第 \(1\) 个分量上为零。这三个自同态之间的关系，已经记录了原来的图表。对一般的有限对象范畴，也可以把全部态射的复合写成一个代数的乘法，使图表成为这个代数上的一个模。

\subsection{复合的线性化与范畴代数}
\label{b4:subsec:0105:01}

\begin{definition}\label{b4:def:category-algebra}
设 \(k\) 为含幺交换环，\(\mathcal C\) 为有有限个对象的小范畴。以 \(\mathcal C\) 的全部态射为基构造自由 \(k\)-模，并在基元素上规定
\[
gf=\begin{cases}
g\circ f,&f:X\to Y,\ g:Y\to Z,\\
0,&f,g\;\text{不能依此次序复合}.
\end{cases}
\]
把乘法双线性地延拓到整个模，所得 \(k\)-代数称为 \(\mathcal C\) 的\term[fanchou-daishu]{范畴代数}[category algebra]，记为 \(k\mathcal C\)。这里每个元素只含有限多个态射的线性组合；并不要求 \(\mathcal C\) 只有有限多个态射。
\end{definition}
\symidx{\(k\mathcal C\)：范畴代数}

\begin{proposition}\label{b4:prop:category-algebra-unit}
设 \(k\) 为含幺交换环，\(\mathcal C\) 为有限对象的小范畴。上述乘法使 \(k\mathcal C\) 成为结合含幺代数，其单位为
\[
1=\sum_{X\in\operatorname{Ob}\mathcal C}e_X,
\qquad e_X=1_X.
\]
各 \(e_X\) 两两正交且幂等。
\end{proposition}
\begin{proof}
对三个基态射 \(f,g,h\)，若它们依次可复合，则 \(h(gf)=(hg)f\) 正是范畴中的结合律。若 \(f,g\) 不能依此次序复合，则 \(f\) 的目标不等于 \(g\) 的源；即使 \(hg\) 可复合，其源仍是 \(g\) 的源，所以 \((hg)f=0\)，两边都为零。若 \(g,h\) 不能依此次序复合，则 \(g\) 的目标不等于 \(h\) 的源；即使 \(gf\) 可复合，其目标仍是 \(g\) 的目标，所以 \(h(gf)=0\)，两边也都为零。双线性延拓遂给出全部元素的结合律。

若 \(f:X\to Y\)，则 \(e_Yf=f=fe_X\)，而其余对象的恒等态射在相应一侧与 \(f\) 相乘均为零。因此所示有限和在两侧都是单位。恒等态射的复合还给出 \(e_X^2=e_X\) 及 \(e_Ye_X=0\)（\(X\ne Y\)）。若范畴为空，所得代数为零代数，单位为 \(0\)。
\end{proof}

有限对象条件用在单位的构造上。若对象无限，则形式和 \(\sum_Xe_X\) 不属于这个只允许有限支集的自由模，不能把它当作代数中的一个元素。另一方面，只有一个对象的群或幺半群，即使有无限多个态射，也仍然给出含幺范畴代数。

\subsection{图表与左模的对应}
\label{b4:subsec:0105:02}

一个协变函子 \(F:\mathcal C\to k\text{-}\mathbf{Mod}\) 给每个对象 \(X\) 配置一个 \(k\)-模 \(F(X)\)，给每个态射 \(f:X\to Y\) 配置一个线性映射 \(F(f):F(X)\to F(Y)\)，并使这些映射服从 \(\mathcal C\) 的复合关系。这样的函子就是一个以 \(\mathcal C\) 为形状的模图表。图表之间的自然变换则是一族与所有图表箭头相容的线性映射。

\begin{theorem}\label{b4:thm:category-algebra-representations}
设 \(k\) 为含幺交换环，\(\mathcal C\) 为有有限个对象的小范畴。以自然变换为态射的函子范畴与含幺左模范畴之间有等价
\[
\operatorname{Fun}(\mathcal C,k\text{-}\mathbf{Mod})
\simeq k\mathcal C\text{-}\mathbf{Mod}.
\]
这个等价把 \(F\) 送到 \(\bigoplus_XF(X)\)，把左 \(k\mathcal C\)-模 \(M\) 送到图表 \(X\mapsto e_XM\)。
\end{theorem}
\begin{proof}
对函子 \(F\)，令 \(M_F=\bigoplus_XF(X)\)。基态射 \(f:X\to Y\) 在 \(X\) 分量上按 \(F(f)\) 作用，所得向量放入 \(Y\) 分量；在其余分量上作用为零。线性延拓定义整个 \(k\mathcal C\) 的作用。若 \(g\) 与 \(f\) 可复合，函子公理给出 \(g(fv)=(gf)v\)；若不可复合，\(fv\) 所在的分量被 \(g\) 送到零。因此作用满足结合律，\(\sum_Xe_X\) 则在各分量上为恒等，故这是含幺左模。

给定自然变换 \(\alpha:F\Rightarrow G\)，映射 \(\bigoplus_X\alpha_X:M_F\to M_G\) 为 \(k\mathcal C\)-线性，因为对每个 \(f:X\to Y\)，自然性恰是
\[
\alpha_YF(f)=G(f)\alpha_X.
\]
直和映射保持恒等与复合，于是得到一个从图表到左模的函子。

反过来，设 \(M\) 为含幺左 \(k\mathcal C\)-模。正交幂等元给出
\[
M=\bigoplus_Xe_XM.
\]
事实上，对象只有有限多个，由 \(1=\sum_Xe_X\) 得到有限和 \(m=\sum_Xe_Xm\)；若 \(\sum_Xm_X=0\) 且 \(m_X\in e_XM\)，乘以 \(e_Y\) 就得到 \(m_Y=0\)。对 \(f:X\to Y\)，定义
\[
F_M(f):e_XM\longrightarrow e_YM,
\qquad m\longmapsto fm.
\]
等式 \(e_Yf=f=fe_X\) 保证值落在指定分量，代数的乘法保证复合相容，\(e_X\) 在 \(e_XM\) 上是恒等。这给出协变函子。模同态 \(h:M\to N\) 满足 \(h(e_Xm)=e_Xh(m)\)，所以可限制到各分量，并与上述箭头相容；因而得到一个从 \(k\mathcal C\) 模范畴回到图表范畴的协变函子。

从 \(F\) 出发再取分量，得到的 \(e_XM_F\) 就是原来的 \(F(X)\)，其箭头与自然变换也相同。从 \(M\) 出发再取直和，则求和映射
\[
\bigoplus_Xe_XM\longrightarrow M,
\qquad (m_X)_X\longmapsto\sum_Xm_X
\]
是模同构，其逆为 \(m\mapsto(e_Xm)_X\)。因对象有限，这族分量自动具有有限支集，确实属于上述直和。这两种识别都与原来的态射相容，故给出拟逆及所需自然同构。
\end{proof}

这里采用左模，是因为态射 \(f:X\to Y\) 应把 \(X\) 分量送入 \(Y\) 分量，而乘法 \(gf\) 表示先作 \(f\) 再作 \(g\)。这一对应的有限对象形式可对照 Webb 的范畴代数构造\cite[Section 2, Proposition 2.1]{Webb2007CategoryRepresentations}。

\subsection{群、路径与偏序的例子}
\label{b4:subsec:0105:03}

把群 \(G\) 看成单对象范畴 \(\mathcal BG\)，便有 \(k\mathcal BG=kG\)：基为群元素，乘法为群乘法。因此，群在 \(k\)-模上的线性表示就是 \(\mathcal BG\) 形状的模图表。自然变换的唯一分量与每个群元素的作用相容，正是表示之间的同态。这里没有要求 \(G\) 有限。
\symidx{\(kG\)：群代数}

\ProseParagraphBreak
再取一个有有限多个顶点的有向图 \(Q\)。以顶点为对象，以有向路径为态射，空路径为恒等，路径的连接为复合，得到自由路径范畴 \(\mathcal P(Q)\)。其范畴代数称为\term[lujing-daishu]{路径代数}[path algebra]，记为 \(kQ\)。从 \(\mathcal P(Q)\) 到 \(k\text{-}\mathbf{Mod}\) 的函子由各顶点的模和各箭头的线性映射唯一决定：任意路径的映射是沿路径依次复合，空路径映成恒等。因而 \(kQ\)-模等价于这个有向图的线性表示。图若有有向环，路径通常有无限多个，但顶点仍有限，所以路径代数仍有单位。
\symidx{\(\mathcal P(Q),\ kQ\)：有向图的路径范畴与路径代数}

\ProseParagraphBreak
对于单箭头 \(0\xrightarrow{a}1\)，基只有 \(e_0,e_1,a\)。对应
\[
e_0\longmapsto E_{11},\qquad
e_1\longmapsto E_{22},\qquad
a\longmapsto E_{21}
\]
把路径代数识别为 \(2\times2\) 下三角矩阵代数。这里 \(E_{ij}\) 表示第 \(i\) 行第 \(j\) 列为 \(1\)、其余为零的矩阵。左模中的两个幂等分量给出 \(V_0,V_1\)，而 \(a\) 给出 \(u:V_0\to V_1\)，恢复节首的图表。若只有一个顶点与一个环箭头 \(a\)，则每条路径唯一写成 \(1,a,a^2,\ldots\)，其中 \(1\) 是空路径。路径乘法满足 \(a^ia^j=a^{i+j}\)，所以 \(a\mapsto t\) 把路径代数识别为 \(k[t]\)；只有这一个生成箭头，并不会出现不同箭头次序的非交换问题。一个左模就是一个 \(k\)-模连同一个指定自同态。

\ProseParagraphBreak
设 \(P\) 为有限偏序集，把它视为偏序范畴。对 \(x\le y\)，记唯一态射为 \(a_{xy}:x\to y\)。于是 \(kP\) 的基由全部这样的区间指标给出，并满足
\[
a_{yz}a_{xy}=a_{xz},
\]
其余不能复合的基元素乘积为零。把偏序线性排列，使 \(x<y\) 时 \(x\) 排在 \(y\) 前，映射 \(a_{xy}\mapsto E_{yx}\) 把 \(kP\) 写成一个下三角矩阵子代数。这是有限偏序集的\term[guanlian-daishu]{关联代数}[incidence algebra]的一种复合约定：矩阵的列指标为源、行指标为目标。若采用把区间 \((x,y)\) 放在第 \(x\) 行、第 \(y\) 列的通常卷积约定，则对应的是此处代数的相反代数。
\symidx{\(kP\)：有限偏序集的关联代数}

\ProseParagraphBreak
偏序图表与自由路径图表之间还有实质区别。在菱形偏序 \(0<1<3\)、\(0<2<3\) 中，从 \(0\) 到 \(3\) 只有一个态射，故两条复合路径必须给出相同线性映射；在只有这四条箭头的自由路径范畴中，两条路径是不同态射，没有这一等式。相应地，前者的范畴代数是后者的路径代数模去“两条路径之差”所得的商。这样，范畴中的复合关系既能约束图表，也能写成代数中的关系。
